% Section 4: Method
% Prose source of truth: ../../drafts/Paper_Method_draft_EN.md
% Two-act structure tied to the Intro thesis:
%   §4.1-4.3 consolidate the signal (objective / buffer / U-shape reweighting)
%   §4.4-4.5 produce a clean signal (sandbox winner-sync / simulated-user data)
% Figures 1-5 live in ../assets/ (see ../assets/README.md for the placeholder list).

\section{Method}
\label{sec:method}

\TODO{port from ../../drafts/Paper\_Method\_draft\_EN.md (keep subsection transitions)}

\subsection{Continual-Learning Objective}
\label{sec:obj}
% L_cl = lambda1 L_rl + lambda2 L_kl + lambda3 L_replay + lambda4 L_ent.
% Drop L_reg; entropy always on. Injected via set_loss_fn, no fork.

\subsection{Capability-Partitioned Prioritized Replay Buffer}
\label{sec:buffer}
% Partition by capability not difficulty; prioritize by anti-forgetting value
% not reward; bucket-local eviction; two-level sampling.

\subsection{U-Shaped Block Reweighting for Replay}
\label{sec:reweight}
% Action-block segmentation; U shape (first + last blocks high); gamma=delta=1
% recovers uniform; clip the product.

\subsection{Sandboxed Rollout with Winner-Synchronized Sessions}
\label{sec:rollout}
% Inference outside / actions inside the sandbox; 16x8 + per-query winner sync.
% Implementation note: replacing the framework's default rollout via the
% official rollout-manager injection point (non-invasive), to realize
% winner-sync inside a standard framework.

\subsection{Online Multi-Turn Query Construction via Simulated Users}
\label{sec:usersim}

The second source of signal corruption is multi-turn data. A follow-up query
typically references the result of the previous turn (e.g., ``the figures on
page 3 of that deck are wrong''). When such a follow-up is fixed at
data-collection time but the realized state $e_k$ is a stochastic function of
the policy, the follow-up's premise holds only with probability
$\Pr[\phi(q_{k+1})(e_k)] < 1$, and this probability \emph{drifts} as the policy
improves. A failed premise produces a corrupted signal: the model is either
rewarded for hallucinating a fix to a non-existent problem, or penalized for
honestly reporting that no problem exists.

We eliminate premise drift by \textbf{constructing follow-ups online, after
observing the realized state}. From each real session we retain only the first
query as a seed (preserving the true distribution of user intents) and generate
subsequent turns at the winner-synchronization boundary with three cooperating
agents. The real sessions \emph{do} contain genuine follow-ups
$q_2,\dots,q_K$, but these were written, at collection time, against the
\emph{original} session's execution result; during training rollout our policy's
response to $q_1$ differs from that original, so the premises of these real
follow-ups fail just as much. Reusing them would re-introduce premise
drift---\textbf{which is exactly why we discard $q_2,\dots,q_K$ and keep only the
seed $q_1$}: the value of the real data lies in the genuine intent distribution
carried by its first query, not in subsequent turns bound to one particular
execution. The three agents are:

\begin{itemize}
  \item an \textbf{observer} (no persona) that compares the winner sandbox
    \textbf{before vs.\ after} the turn and gathers a neutral, structured report
    $R_t = \mathrm{Obs}(e_t^w)$ from the resulting environment \textbf{diff}
    (files created/modified/removed with content, plus system-state changes;
    read-only), capturing intermediate as well as final results. Crucially the
    observer is grounded on the \textbf{real state only}: it does \textbf{not}
    receive the actor's trajectory or claims at all. This makes anti
    reward-hacking \emph{structural}---since the actor's narrative never enters
    the observer's judgment, completion can only reflect the verified effect, so
    a policy cannot earn reward by \emph{claiming} completion it did not deliver.
    (The trajectory is still made available to the reward judge for
    safety/robustness, carried pass-through and never shown to the observer
    model.) The observer is also independent of the policy model (different
    model/backend), decoupling observation from action;
  \item a \textbf{questioner} that carries one of $42$ fixed personas
    (profession, preference, user profile, and an \emph{observation focus}: whole
    vs.\ detail, form vs.\ content). Exactly one persona is drawn at the start of
    each session and held fixed throughout (one persona per session, $p$ constant
    within a session); it produces the next query
    $q_{t+1}\sim Q(\cdot\mid p, R_t, H_t)$, where $H_t$ is the winner-derived
    session history; and
  \item a \textbf{reward model} that scores the turn from the same report over
    two channels: the observer's state diff (ground truth for \emph{completion})
    and the actor trajectory carried pass-through on $R_t$ (for
    \emph{safety/robustness}), $r_t = \mathrm{Reward}(R_t, \text{rubric})$. A turn
    that changed nothing (empty diff) is short-circuited to zero without a judge
    call.
\end{itemize}

The report $R_t$ serves both downstream consumers, guaranteeing that the facts
used to score a turn and the facts used to pose the next question are identical,
and decoupling objective evidence-gathering (observer) from subjective stance
(questioner). Because the follow-up is written after the observer has inspected
$e_t^w$, its premise is grounded by construction,
$\Pr[\phi(q_{t+1})(e_t^w)]\approx 1$. Grounding the reward's \emph{completion} in
the observed effect (whether a file was actually produced, whether its contents
are correct) rather than any textual claim is what guards against reward
hacking---and because the observer itself never reads the trajectory, that
grounding is structural rather than a matter of the judge's diligence.

\paragraph{Cognitive intuition: eyes and brain.}
This division of labor mirrors the human \textbf{eye} and \textbf{brain}. The
observer is the \textbf{eye}: objective, unemotional, stating only facts---it
does not listen to the actor's narration but performs a deterministic diff of the
sandbox environment, faithfully imaging \emph{what actually happened} (files
added/modified with content, system-state changes, claim-vs-reality
discrepancies), like a retina that registers light without adding preference. The
questioner is the \textbf{brain}: facing the same objective image, it interprets
subjectively through the lens of a persona---the persona sets the focus (a
detail-oriented auditor fixates on a number, a big-picture manager reacts to the
overall deliverable), the standard (a researcher demands the raw file, an SRE
asks whether the fix actually runs), and the emotional register (calm personas
coach, hot-tempered ones press). \textbf{The eye guarantees what is seen is true;
the brain guarantees the judgment behaves like a real person}: the former makes
anti-reward-hacking structural (judgment anchored to ground truth, not
narration), the latter makes follow-ups diverse and faithful to the real online
distribution. One answers ``what is,'' the other ``how one sees it.''

\paragraph{Four-dimensional anti-collapse.}
Self-dialogue is known to suffer from mode collapse---follow-ups converge to a few
templated patterns, and entropy decays with turn count. We counter this with four
orthogonal mechanisms:
\begin{enumerate}
  \item \textbf{Persona diversity}: 42 personas (differing in profession,
    preference, and observation focus) are drawn one per session, causing the
    same objective report to be queried from different angles;
  \item \textbf{Turn budget}: session length is fully controlled by the
    questioner---when satisfied with the observed output it issues
    \texttt{<end\_session>} naturally; when unsatisfied it keeps asking. No fixed
    $K$ budget, no random truncation: the persona decides, not a hard cap;
  \item \textbf{State evolution}: each turn's winner state has been altered by the
    previous follow-up, so even a fixed persona sees a different report each turn;
  \item \textbf{Questioner multi-model rotation}: the questioner rotates across 4
    cross-vendor models (\texttt{claude-4.6-sonnet}, \texttt{deepseek-v4-pro},
    \texttt{qwen3-max}, \texttt{kimi-k2.6}), switching every 5 queries.
\end{enumerate}
The first three dimensions have precedent in the literature (persona
diversity~\cite{A2}, turn truncation~\cite{B4}, input drift). The
fourth---model-level heterogeneity---is our addition: a single LLM, even at
temperature $0.9$, has a fixed output-distribution center and oscillates within
one semantic space. Models from different vendors differ in training data,
alignment, and language style, so rotating across them disperses the
\emph{focus, tone, and angle} of follow-ups in a way temperature alone cannot.
Personas determine \emph{what to ask}; models determine \emph{how to ask it}---the
two inject diversity along orthogonal axes, and their joint effect exceeds either
alone. The 5-query cadence balances coherence (one model's consecutive queries
preserve local context) against heterogeneity (the periodic switch prevents any
single model from dominating).

\paragraph{Patience-governed handling of failed turns.}
Rather than a fixed retry cap, we make the decision to retry or abandon a failed
turn an attribute of the simulated user, governed by two per-persona quantities:
initial patience $P_0(p) \in [3, 10]$ and a retention rate $r(p) \in
(0.45, 0.90)$. Let $k=1,2,\dots$ index consecutive failed turns within a session.
On the $k$-th failure, patience follows a multiplicative decay,
\begin{equation}
  P_k = P_0(p) \cdot r(p)^{\,k}.
\end{equation}
If $P_k \le 0.1$ the session terminates deterministically; if $P_k \ge 1$ the
controller always retries; otherwise it retries with probability $P_k$, and the
questioner generates a persona-voiced natural retry (an impatient persona snaps
``又挂了？再来'' while a calm one says ``没关系，再试一次''). Making \emph{both}
$P_0$ and $r$ persona-specific decouples two distinct traits---\emph{initial
tolerance} (high $P_0$ grants more chances) and \emph{retention} ($r$ near $1$
sustains patience, $r$ near $0.5$ exhausts it quickly). Multiplicative decay
captures cumulative fatigue: each failure reduces the \emph{remaining} patience
proportionally rather than subtracting a fixed amount, aligning with the
``dwindling goodwill'' intuition and giving a clean closed form. Each session
independently loads its $P_0$ from the persona; patience never carries over across
sessions. Successful turns end the session via \texttt{<end\_session>} without
touching the patience mechanism.

A useful side effect is an \textbf{adaptive curriculum}: the questioner always
critiques the \emph{current} policy's actual output, so as the policy strengthens
and residual flaws become subtler, the follow-ups become correspondingly
finer---tracking the capability frontier without a hand-designed difficulty
schedule.


